% Generated draft: requires mathematical and source review.
\documentclass[11pt]{article}
\usepackage[margin=25mm]{geometry}
\usepackage{fontspec}
\setmainfont{Latin Modern Roman}
\usepackage{amsmath,amssymb,mathtools,unicode-math}
\setmathfont{Latin Modern Math}
\usepackage{enumitem,needspace,xurl,hyperref}
\setlength{\parindent}{0pt}
\setlength{\parskip}{5pt}
\setlength{\emergencystretch}{3em}
\newcommand{\N}{\mathbb{N}}
\newcommand{\Z}{\mathbb{Z}}
\newcommand{\Q}{\mathbb{Q}}
\newcommand{\R}{\mathbb{R}}
\newcommand{\C}{\mathbb{C}}
\newcommand{\F}{\mathbb{F}}
\newcommand{\A}{\mathbb{A}}
\newcommand{\PP}{\mathbb{P}}
\newcommand{\E}{\mathbb{E}}
\newcommand{\eps}{\varepsilon}
\newcommand{\dd}{\,\mathrm{d}}
\begin{document}
\section*{UnsolvedMath ID: 30005325}
\textbf{Source problem code:} catalogue entry
\subsection{Definitions and mathematical statement}
Let \(X=[0,1]^2\) carry its usual topology and Borel \(\sigma\)-algebra, and let \(\lambda\) denote Lebesgue probability measure on \([0,1]\). A \emph{permuton} is a Borel probability measure \(\mu\) on \(X\) whose two coordinate marginals are both \(\lambda\); explicitly, for every Borel set \(A\subseteq[0,1]\),
\[
\mu(A\times[0,1])=\lambda(A),\qquad
\mu([0,1]\times A)=\lambda(A).
\]
Its topological support is
\[
\operatorname{supp}(\mu)=\{z\in X:\mu(B_{\R^2}(z,r)\cap X)>0\text{ for every }r>0\}.
\]
\textbf{Problem.} Characterize all closed sets \(S\subseteq X\) for which there exists a permuton \(\mu\) satisfying
\[
\operatorname{supp}(\mu)=S.
\]
Thus the requested measure must satisfy both \(\mu(X\setminus S)=0\) and \(\mu(B_{\R^2}(s,r)\cap X)>0\) for every \(s\in S\) and every \(r>0\).

This formulation follows the supplied clean statement, which asks for exact topological support in the square. The supplied original excerpt instead says a closed set in \([0,1]\), while its figure caption refers to subsets of the unit square. That dimensional discrepancy remains a source issue. The phrase ``support a permuton'' in the caption could also mean merely carrying a permuton, whereas exact support is explicitly required by the clean statement. The broader introductory discussion of arbitrary given real distributions is treated as motivation; the stated permuton problem requires uniform marginals.
\subsection{Short English statement}
Determine precisely which closed subsets of the unit square are the exact supports of probability measures having uniform marginals on both coordinates. The supplied clean statement specifies the square, although the original excerpt contains a conflicting interval reference.
\subsection{Sources}
\begin{itemize}
\item Local UnsolvedMath registry, supplied record with permanent ID \(30005325\), title \emph{Closed Sets Supporting Permutons}, source problem code \texttt{catalogue entry}. The registry supplies the status \texttt{partially\_solved} and the statement label \texttt{corrected\_verified}; these are supplied curation labels, not independent verification here.
\item Supplied report citation: \emph{Enumerative Combinatorics} (2023), \emph{Oberwolfach Reports}, pp.\ 3215--3305, DOI \texttt{10.4171/owr/2022/57}. Original report link: \url{https://doi.org/10.4171/owr/2022/57}. The supplied excerpt attributes \emph{Sets that Support a Joint Distribution} to Peter Winkler, joint work with Chris Coscia and Martin Tassy. The record gives proposed year 2022 and no separate proposer entry.
\item Supplied literature source: Coscia--Tassy--Winkler, \url{https://arxiv.org/abs/2608.04135}. No paper title or theorem identifiers are supplied.
\item Literature assessment dated exactly \texttt{2026-08-22}: the supplied curation reports necessary and sufficient conditions in the discrete setting and a partial characterization for continuous supports, but no complete characterization of all closed subsets of the square. This is supplied dated curation, not independently verified. The background explicitly calls it dated literature triage derived from the supplied OWR curation and requests independent verification of time-sensitive status and cited claims.
\end{itemize}
\end{document}
